\documentclass[11pt,a4paper]{article}
\usepackage[margin=27mm,headheight=15pt]{geometry}
\usepackage[T1]{fontenc}
\usepackage[utf8]{inputenc}
\usepackage{lmodern}
\usepackage{amsmath,amssymb,amsthm,mathtools}
\usepackage{booktabs,array,longtable}
\usepackage{microtype}
\usepackage{listings}
\usepackage{xurl}
\usepackage[hidelinks]{hyperref}
\usepackage{fancyhdr}
\usepackage{enumitem}
\pagestyle{fancy}
\fancyhf{}
\fancyhead[L]{\small OEIS A352969}
\fancyhead[R]{\small Correctness and quantitative growth}
\fancyfoot[C]{\thepage}
\setlength{\parskip}{0.35em}
\setlength{\parindent}{1.25em}
\setlength{\emergencystretch}{2em}
\allowdisplaybreaks[1]
\raggedbottom
\newtheorem{theorem}{Theorem}[section]
\newtheorem{lemma}[theorem]{Lemma}
\newtheorem{proposition}[theorem]{Proposition}
\newtheorem{corollary}[theorem]{Corollary}
\theoremstyle{definition}
\newtheorem{definition}[theorem]{Definition}
\theoremstyle{remark}
\newtheorem{remark}[theorem]{Remark}
\newcommand{\NN}{\mathbb{N}}
\renewcommand{\lg}{\log_2}
\newcommand{\card}[1]{\lvert #1\rvert}
\newcommand{\ceil}[1]{\left\lceil #1\right\rceil}
\newcommand{\floor}[1]{\left\lfloor #1\right\rfloor}
\newcommand{\D}{\mathcal D}
\DeclareMathOperator{\height}{ht}
\lstset{language=Python,basicstyle=\small\ttfamily,columns=fullflexible,
  keepspaces=true,showstringspaces=false,breaklines=true,frame=single,
  framerule=0.3pt,xleftmargin=0.6em,xrightmargin=0.6em,
  aboveskip=1em,belowskip=0.8em,tabsize=4}
\hypersetup{pdftitle={Correctness and Quantitative Growth of OEIS A352969},
 pdfauthor={Research note prepared for Vladimir Reshetnikov},
 pdfsubject={Iterated pairwise sums and products; constructive formula depth; asymptotics}}

\title{\textbf{Iterated Sums and Products from One}\\[0.45em]
  \large Correctness and Quantitative Growth of OEIS A352969}
\author{A research note prepared for Vladimir Reshetnikov}
\date{September 20, 2026}

\begin{document}
\maketitle
\begin{abstract}
Let $S_0=\{1\}$ and let $S_{n+1}$ consist of all sums and products of two
members of $S_n$, allowing equal operands. We prove that the Python program
linked from OEIS A352969 computes $S_n$ exactly, under its natural input
and cache-integrity assumptions. We then establish unconditional quantitative
growth estimates. Writing $a_n=\card{S_n}$, we prove
\[
 \lg\lg a_n=n+O(\sqrt{\log n}),\qquad
 \frac{\log a_{n+1}}{\log a_n}\longrightarrow 2.
\]
The first estimate is made explicit: for $n\ge16$, if
$t=\ceil{\lg(2n)}$ and $r$ is the smallest integer with $t\le r(r+1)/2$,
then $2^{\,2^{n-r-4}}\le a_n\le2^{\,2^{n-1}}$.
The proof combines a low-depth radix construction for arbitrary positive
integers with collision-free products of prime multisets. An elementary
divisor bound then controls the loss in one multiplication step.
The normalized logarithms $2^{-n}\lg a_n$ decrease to a limit, but its
value and strict positivity are not determined here. Thus the result is a
sharp logarithmic-scale growth theorem, not a claimed multiplicative
asymptotic equivalent for $a_n$. Reproducible code, independently checked
small-stage data, and an executable expression certificate accompany the article.
\end{abstract}

\clearpage
\tableofcontents
\clearpage

\section{The sequence and the main results}\label{sec:main}
For a finite set $A$ of positive integers, write
\[
 A+A=\{x+y:x,y\in A\},\qquad
 AA=\{xy:x,y\in A\},\qquad \Phi(A)=(A+A)\cup AA.
\]
The sequence under consideration is
\begin{equation}\label{eq:definition}
 S_0=\{1\},\qquad S_{n+1}=\Phi(S_n),\qquad a_n=\card{S_n}.
\end{equation}
This is OEIS A352969, entered by Vladimir Reshetnikov in April 2022.
The entry attributes its Python program to Chai Wah Wu and currently lists
\begin{equation}\label{eq:terms}
 1,\ 2,\ 4,\ 11,\ 52,\ 678,\ 67144,\ 357306081
\end{equation}
for $n=0,\ldots,7$~\cite{oeis}. We independently recomputed the first seven
of these values, through $n=6$; the value at $n=7$ is used only as reported
OEIS data. All logarithms denoted by $\log$ are natural, and $\lg$ denotes
base two.

Two distinct tasks must be separated. Proving the program correct amounts
to checking that its enumeration implements \eqref{eq:definition} without
omissions or spurious values. Estimating $a_n$ is harder: the program
counts \emph{values}, not arithmetic expressions, and many expressions
have the same value.

\begin{theorem}[Quantitative growth]\label{thm:main}
For $n\ge16$, define
\[
 t_n=\ceil{\lg(2n)},\qquad
 r_n=\min\{r\in\NN:t_n\le r(r+1)/2\}.
\]
Then
\begin{equation}\label{eq:explicit-main}
 \boxed{\quad 2^{\,2^{n-r_n-4}}\le a_n\le2^{\,2^{n-1}}.\quad}
\end{equation}
In particular,
\begin{align}
 n-\sqrt{2\lg n}-O(1)&\le \lg\lg a_n\le n-1,\label{eq:doublelog}\\
 \lim_{n\to\infty}\frac{\log\log a_n}{n}&=\log2,\label{eq:limit-main}\\
 \log a_n&=\frac{2^n}{n^{o(1)}}.\label{eq:subpoly}
\end{align}
Moreover, there is a nonnegative sequence $\varepsilon_n$ such that
\begin{equation}\label{eq:local-main}
 a_{n+1}=a_n^{\,2-\varepsilon_n},\qquad
 0\le\varepsilon_n\le n^{-1+o(1)}.
\end{equation}
Consequently $\log a_{n+1}/\log a_n\to2$.
\end{theorem}

The notation in \eqref{eq:subpoly} means that the positive factor
$2^n/\log a_n$ grows more slowly than every fixed positive power of $n$;
it is not an assertion that this factor tends to infinity.
An equivalent and less compressed statement is that for every $\delta>0$
and all sufficiently large $n$,
\[
 \frac{2^n}{n^\delta}\le\log a_n\le\frac{\log2}{2}\,2^n.
\]
The threshold in the lower bound may depend on $\delta$.

\begin{theorem}[Normalized logarithms]\label{thm:entropy-intro}
The sequence
\[
 \beta_n=2^{-n}\lg a_n\qquad(n\ge1)
\]
is nonincreasing and has a limit $\theta\ge0$. Thus
\begin{equation}\label{eq:theta-form}
 a_n=2^{(\theta+o(1))2^n}.
\end{equation}
The independently recomputed data imply
$\theta\le0.2505464198$; the reported OEIS value of $a_7$ improves this to
$\theta\le0.2219733221$.
\end{theorem}

Formula \eqref{eq:theta-form} does not by itself prove $\theta>0$.
Even a positive value of $\theta$ would give an equivalent for $\log a_n$,
not automatically a multiplicative equivalent
$a_n\sim2^{\theta2^n}$. These distinctions matter for a doubly exponential
sequence.

\section{What the Python program does}\label{sec:program}
The OEIS listing defines a cached function named \texttt{A352969\_set}.
Its base case returns $\{1\}$. At a positive argument it forms two sets:
one from the sums, the other from the products, of
\texttt{combinations\_with\_replacement(previous, 2)}, and returns their
union. A final wrapper takes the cardinality. In the retrieved listing
that wrapper is named \texttt{A353969}; this is a naming typo, not a
mathematical change~\cite{oeis}.

The following independently written version expresses the same recurrence
in one pass over the pairs. It deliberately returns immutable results and
validates the public input before entering the cache.
\begin{lstlisting}
from functools import lru_cache
from itertools import combinations_with_replacement

@lru_cache(maxsize=None)
def _stage(n):
    if n == 0:
        return frozenset((1,))
    result = set()
    for x, y in combinations_with_replacement(_stage(n - 1), 2):
        result.add(x + y)
        result.add(x * y)
    return frozenset(result)

def reachable(n):
    if isinstance(n, bool) or not isinstance(n, int):
        raise TypeError("n must be an integer")
    if n < 0:
        raise ValueError("n must be nonnegative")
    return _stage(n)

def A352969(n):
    return len(reachable(n))
\end{lstlisting}
The package contains both cached and iterative variants, with additional
input checks and an optional pair-count budget. The compact listing above
is for exposition; the executable implementation is
\texttt{code/a352969.py}.

\subsection{Why unordered pairs are sufficient}
Fix an arbitrary enumeration $A=\{u_1,\ldots,u_m\}$.
The iterator with replacement and length two visits exactly the index pairs
\[
 (i,j)\quad\text{with}\quad1\le i\le j\le m.
\]
Its ordering is determined by the input iterable, not by the numerical
order of its elements~\cite{itertools}. A Python set therefore need not be
sorted before it is passed to the iterator.

For every ordered pair $(x,y)\in A^2$, exactly one of $(x,y)$ and $(y,x)$
occurs in this unordered enumeration, except that $(x,x)$ occurs once.
Both operations commute, so replacing an ordered pair by its reversal
changes neither output. Equal pairs must be retained: for example, the
first step needs $1+1$ and $1\cdot1$.

\begin{lemma}[One-step enumeration]\label{lem:onestep}
For any finite set $A$ of positive integers, the union of the sum outputs
and product outputs from the unordered enumeration with replacement is
exactly $\Phi(A)$.
\end{lemma}
\begin{proof}
Every emitted sum or product has both operands in $A$, so no emitted value
lies outside $\Phi(A)$. Conversely, if $z=x+y$ or $z=xy$ with $x,y\in A$,
the iterator visits the unordered pair containing $x$ and $y$.
Commutativity ensures that the corresponding output equals $z$.
Set insertion discards repeated values but never discards their membership.
Taking the union retains values produced by either operation.
\end{proof}

\subsection{Correctness and termination}
\begin{theorem}[Correctness of the OEIS algorithm]\label{thm:correct}
For every nonnegative integer $n$, the OEIS set function returns $S_n$,
provided its previously cached sets have not been externally modified and
the execution has sufficient memory and stack capacity. Its cardinality
wrapper therefore returns $a_n$.
\end{theorem}
\begin{proof}
Induct on $n$. At $n=0$, the returned set is $\{1\}=S_0$.
For $n>0$, the induction hypothesis identifies the recursively returned
set with $S_{n-1}$. Lemma~\ref{lem:onestep} identifies the two comprehensions
and their union with $\Phi(S_{n-1})=S_n$.
The wrapper applies cardinality to this exact set.

For mathematical termination, the recursion decreases the nonnegative
integer argument by one. At each level the previous set is finite, and
there are only $m(m+1)/2$ unordered pairs. Thus each loop is finite.
Python's integer operations are exact on integer operands; there is no
fixed-width overflow that silently changes the values~\cite{types,math}.
Resource exceptions are a separate issue from the mathematical termination
of this finite procedure.
\end{proof}

The cache changes reuse, not the recurrence. In a single-threaded cold
call, each stage is actually computed once; the second reference to the
previous stage in the original two-comprehension implementation reuses its
cached value. The cache retains references to returned objects rather
than defensive copies~\cite{functools}. This explains the qualification
about external mutation: inserting a foreign integer into a cached mutable
set invalidates the induction hypothesis for later calls. Returning a
\texttt{frozenset} removes that particular failure mode~\cite{types}.

The original listing is intended for $n\in\NN$. A negative input never
reaches the base case by decrementing, and a noninteger input is not part
of its mathematical contract. The public validation in the supplied
implementation makes this contract explicit. The distinction between
booleans and integers is checked deliberately, since Python booleans are
an integer subtype~\cite{types}.

\section{Structural facts about the generated sets}\label{sec:structure}
\subsection{Persistence and arithmetic-expression height}
\begin{lemma}\label{lem:nested}
For every $n$, $1\in S_n$ and $S_n\subseteq S_{n+1}$.
\end{lemma}
\begin{proof}
The first assertion follows from $1\cdot1=1$.
If $x\in S_n$, then $x=x\cdot1\in S_{n+1}$.
\end{proof}

A legal expression has leaves all equal to $1$ and internal vertices labeled
$+$ or $\times$, each with two children. Its height is zero for a leaf
and one plus the maximum child height at an internal vertex. Define
\[
 h(m)=\min\{\height(E):E\text{ is a legal expression with value }m\}.
\]
Every positive integer has such an expression, for example a sum of ones.
This $h$ minimizes height without first minimizing the number of leaves.
It is therefore not the ordinary integer complexity recorded in
A005245, nor the leaf-optimal depth convention of A210659~\cite{complexity,depth}.

\begin{proposition}\label{prop:height}
For every $n\ge0$,
\begin{equation}\label{eq:height-set}
 S_n=\{m\ge1:h(m)\le n\}.
\end{equation}
For $m>1$, $h(m)$ is also specified by the finite recurrence
\begin{equation}\label{eq:h-dp}
 h(m)=1+\min\left\{
 \begin{array}{ll}
 \max(h(u),h(v)) &: u+v=m,\quad u,v\ge1,\\
 \max(h(u),h(v)) &: uv=m,\quad u,v\ge2.
 \end{array}\right.
\end{equation}
\end{proposition}
\begin{proof}
An expression of height at most $n+1$ has root operands of height at most
$n$, except for the leaf expression, whose value $1$ persists by
Lemma~\ref{lem:nested}. Induction on $n$ proves \eqref{eq:height-set} in
both directions. Alternatively, shorter expressions can be padded by
multiplications by $1$.

An optimal expression for $m>1$ cannot have a root multiplication by $1$:
removing that root would strictly lower its height. Every remaining root
has one of the proper decompositions in \eqref{eq:h-dp}, whose operands
are smaller than $m$. Optimizing over these roots proves the recurrence.
\end{proof}

The program for \eqref{eq:h-dp} supplies a useful independent test oracle:
it indexes the computation by the numerical value $m$, whereas the OEIS
algorithm indexes it by the stage $n$.

\subsection{Maximum value and powers of two}
\begin{proposition}\label{prop:maximum}
Let $M_n=\max S_n$. Then
\begin{equation}\label{eq:max}
 M_0=1,\qquad M_n=2^{2^{n-1}}\quad(n\ge1).
\end{equation}
In particular $a_n\le M_n$. For every integer $e\ge1$,
\begin{equation}\label{eq:powerheight}
 h(2^e)=1+\ceil{\lg e}.
\end{equation}
\end{proposition}
\begin{proof}
The largest sum and product from $S_n$ are $2M_n$ and $M_n^2$, respectively.
Thus $M_{n+1}=\max(2M_n,M_n^2)$. The first step gives $M_1=2$;
thereafter $M_n^2\ge2M_n$, giving \eqref{eq:max} by repeated squaring.
Every member is a positive integer at most $M_n$, so $a_n\le M_n$.

For \eqref{eq:powerheight}, form a balanced product of $e$ copies of
$2=1+1$. It has height at most $1+\ceil{\lg e}$.
Conversely $2^e\in S_j$ requires $2^e\le2^{2^{j-1}}$, hence
$j\ge1+\ceil{\lg e}$. The two bounds coincide.
\end{proof}

\subsection{Cardinality bounds and the entropy limit}
\begin{proposition}\label{prop:cardbounds}
For all $n\ge0$, $a_{n+1}\ge2a_n$. For all $n\ge1$,
\begin{equation}\label{eq:square}
 a_{n+1}\le a_n^2.
\end{equation}
\end{proposition}
\begin{proof}
Write the elements of $A=S_n$ increasingly as
$1=u_1<u_2<\cdots<u_m$. The sums
\[
 u_1+u_1<\cdots<u_1+u_m<u_2+u_m<\cdots<u_m+u_m
\]
are $2m-1$ distinct members of $A+A$. The additional value $1\in AA$ is
not in $A+A$, proving $\card{\Phi(A)}\ge2m$.

For the upper bound, take $n\ge1$, so $2\in A$.
Each of $A+A$ and $AA$ has at most $m(m+1)/2$ members.
Their intersection contains the $m$ distinct values $2x$, $x\in A$:
they occur both as $x+x$ and as $2\cdot x$. Inclusion--exclusion gives
\[
 \card{\Phi(A)}\le m(m+1)-m=m^2.
\]
\end{proof}

Dividing the logarithm of \eqref{eq:square} by $2^{n+1}$ proves that
$\beta_{n+1}\le\beta_n$. Nonnegativity gives its limit $\theta$, proving
the existence assertion of Theorem~\ref{thm:entropy-intro}.
For every fixed $j\ge1$,
\begin{equation}\label{eq:finite-upper}
 a_n\le a_j^{2^{n-j}}\qquad(n\ge j).
\end{equation}
The numerical upper bounds for $\theta$ follow by substituting the data
at $j=6$ or $j=7$.

A useful consequence is that the sets are extremely sparse relative to
the interval allowed by their largest element. Since $a_3=11$ and
$M_3=16$, \eqref{eq:finite-upper} and \eqref{eq:max} give
\begin{equation}\label{eq:density}
 \frac{a_n}{M_n}\le\left(\frac{11}{16}\right)^{2^{n-3}}
 \longrightarrow0\qquad(n\ge3).
\end{equation}
Thus having the same leading double-logarithmic scale as $M_n$ does not
mean that $S_n$ has appreciable density in $[1,M_n]$.

\section{Why a lower bound needs a construction}\label{sec:collision}
The raw number of candidate outputs at a stage is $a_n(a_n+1)$.
This is an upper bound, not a lower bound. Commutativity, associativity,
distributivity, and numerical coincidences create many collisions.
For instance $2+2=2\cdot2$, and $2\cdot6=3\cdot4$.
Neither counting expression trees nor treating candidate values as
independent random integers proves a lower bound for $a_n$.

We will instead use a family whose collisions are completely understood.
The strategy has three parts: construct a long initial interval at a
moderate height, select the primes in that interval, and multiply a
prescribed number of them in a balanced tree. Unique factorization makes
the last counting step exact. In particular, no assumption about the
primality of specially structured expressions is needed: \emph{every}
integer in the chosen interval is constructible.

\begin{lemma}[Prime-multiset amplification]\label{lem:amplify}
If $S_k$ contains $P$ distinct primes, then for every $n\ge k$, putting
$R=2^{n-k}$ gives
\begin{equation}\label{eq:primecount}
 a_n\ge\binom{P+R-1}{R}\ge\left(\frac{P}{R}\right)^R.
\end{equation}
\end{lemma}
\begin{proof}
Choose a multiset of exactly $R$ primes from the $P$ available primes,
and multiply its elements in a complete balanced tree with $R$ leaves.
The multiplication tree adds exactly $n-k$ to the maximum input height,
so its value belongs to $S_n$.

The prime multiplicities are recovered uniquely from the resulting
integer. Hence different multisets give different values. The number of
multisets is the binomial coefficient in \eqref{eq:primecount}. Finally,
\[
 \binom{P+R-1}{R}
 =\frac{P(P+1)\cdots(P+R-1)}{R!}
 \ge\frac{P^R}{R^R}.
\]
\end{proof}

The first inequality remains valid for any fixed family of pairwise
coprime integers greater than one. Primes are convenient because an
initial interval contains sufficiently many of them, as the next section
proves without using the prime number theorem.

\section{An elementary supply of primes}\label{sec:primes}
\begin{lemma}\label{lem:chebyshev}
For every even integer $X\ge16$,
\begin{equation}\label{eq:prime-lower}
 \pi(X)\ge\frac{X}{2\log X},
\end{equation}
where $\pi(X)$ counts primes at most $X$.
\end{lemma}
\begin{proof}
Write $X=2m$. The central binomial coefficient is the largest of the
$2m+1$ binomial coefficients in the expansion of $(1+1)^{2m}$. Therefore
\[
 \binom{2m}{m}\ge\frac{2^{2m}}{2m+1}=\frac{2^X}{X+1}.
\]
For a prime $p$, counting multiples of each prime power in the factorials
gives
\[
 v_p\binom{2m}{m}
 =\sum_{j\ge1}\left(\floor{2m/p^j}-2\floor{m/p^j}\right).
\]
Each summand is either zero or one and vanishes for $p^j>X$.
Consequently $v_p\binom{2m}{m}\le\floor{\log_p X}$, so the full power of
$p$ dividing the binomial coefficient is at most $X$. All its prime
factors are at most $X$, and hence
\[
 \binom{2m}{m}\le X^{\pi(X)}.
\]
Taking logarithms of both estimates yields
\[
 \pi(X)\ge\frac{X\log2-\log(X+1)}{\log X}.
\]
For $X\ge16$, the numerator is at least $X/2$: this holds at $16$, and
$X(\log2-1/2)-\log(X+1)$ is increasing on $[16,\infty)$.
This proves \eqref{eq:prime-lower}.
\end{proof}

Only a positive constant times $X/\log X$ is needed later. The explicit
constant $1/2$ keeps the finite bounds transparent; optimizing it would
not change the principal growth rate.

\section{A first complete growth proof}\label{sec:first-growth}
Before sharpening the construction, it is useful to see why the base of
the double exponential is already forced to be two.

\begin{lemma}[Binary expansion bound]\label{lem:binary}
For every integer $t\ge1$, every positive integer $m<2^{2^t}$ has an
expression of height at most $2t+1$.
\end{lemma}
\begin{proof}
Expand $m$ in binary. It is a sum of at most $2^t$ powers $2^e$ with
$0\le e<2^t$. The power $1=2^0$ has height zero, and all the other powers
have height at most $t+1$ by \eqref{eq:powerheight}.
A balanced sum of at most $2^t$ positive terms adds at most $t$ levels.
Zero binary digits are omitted; zero is not used as an expression leaf.
\end{proof}

Fix $n\ge16$, put $t=\ceil{\lg(2n)}$, and let $k=2t+1$.
All primes at most $X=2^{2n}$ are strictly below $2^{2^t}$, even when
$2^t=2n$, because the upper endpoint $X$ is composite. They therefore
belong to $S_k$. Also $k\le n$, and
\[
 P=\pi(X)\ge\frac{2^{2n}}{4n\log2},\qquad R=2^{n-k}.
\]
Since $2^k\ge8n^2$, one has $P/R\ge2^n$. Lemma~\ref{lem:amplify} now gives
\[
 \lg a_n\ge R\lg(P/R)\ge nR.
\]
Finally $2^t<4n$ implies $2^k=2(2^t)^2<32n^2$, giving
\begin{equation}\label{eq:coarse}
 \boxed{\qquad a_n\ge2^{2^n/(32n)}\qquad(n\ge16).\qquad}
\end{equation}
For completeness, $k\le n$ follows from
$k\le2\lg n+5\le n$ for $n\ge16$.
Combining \eqref{eq:coarse} with \eqref{eq:max} already proves
\[
 n-\lg n-5\le\lg\lg a_n\le n-1,
\]
and hence \eqref{eq:limit-main}. The next section improves the logarithmic
error to a square root of a logarithm. That improvement will also make a
one-step growth theorem possible.

\section{A sharper construction for arbitrary integers}\label{sec:radix}
Define the uniform height requirement
\[
 \D(t)=\max_{1\le m<2^{2^t}}h(m)\qquad(t\ge0).
\]
The interval is finite and nonempty, and $\D(0)=0$.
Define the triangular index
\begin{equation}\label{eq:triangular}
 r(t)=\min\{r\ge0:t\le r(r+1)/2\}
 =\ceil{\frac{\sqrt{8t+1}-1}{2}}.
\end{equation}

\begin{theorem}[Uniform low-depth construction]\label{thm:radix}
For every integer $t\ge0$,
\begin{equation}\label{eq:uniform}
 \D(t)\le t+r(t)+2.
\end{equation}
In particular, every positive integer with at most $2^t$ binary digits
has such a representation using only $1$, addition, and multiplication.
\end{theorem}
\begin{proof}
The case $t=0$ is immediate. Induct on $t\ge1$ and put $r=r(t)$.
Minimality of $r$ gives
\[
 \frac{(r-1)r}{2}<t\le\frac{r(r+1)}2.
\]
In particular $1\le r\le t$. Set $q=t-r$. Then
\begin{equation}\label{eq:qtriangle}
 0\le q\le\frac{(r-1)r}{2},\qquad r(q)\le r-1.
\end{equation}

Split the binary digits of $m<2^{2^t}$ into $2^r$ blocks, each of width
$w=2^q$ bits. Equivalently,
\begin{equation}\label{eq:radixexpand}
 m=\sum_{j=0}^{2^r-1}d_j2^{jw},\qquad0\le d_j<2^w.
\end{equation}
For every nonzero digit $d_j$, the induction hypothesis and
\eqref{eq:qtriangle} give
\[
 h(d_j)\le\D(q)\le q+r(q)+2\le(t-r)+(r-1)+2=t+1.
\]
If $j>0$, then $1\le jw<2^t$, so
$h(2^{jw})\le t+1$ by \eqref{eq:powerheight}.
The product $d_j2^{jw}$ consequently has height at most $t+2$.
The $j=0$ term needs no multiplication, but the same upper bound applies.

Discard zero terms in \eqref{eq:radixexpand}; at least one term remains.
Their number is at most $2^r$, so a balanced sum adds at most $r$ levels.
The complete expression has height at most $t+r+2$, as required.
\end{proof}

The particular triangular numbers are not an unexplained guess. The
construction uses $r$ levels to add $2^r$ blocks. It is viable whenever
the smaller coefficient problem at parameter $t-r$ has an overhead one
unit below the current overhead. The recurrence
\[
 T_r=T_{r-1}+r,\qquad T_0=0,
\]
for the largest permitted parameter gives exactly $T_r=r(r+1)/2$.

\subsection{What the construction proves about individual integers}
For $m\ge2$, take $t=\ceil{\lg(\lg(m+1))}$, or equivalently choose the
least $t$ for which $m<2^{2^t}$. The maximum-value obstruction and
Theorem~\ref{thm:radix} give
\begin{equation}\label{eq:height-asymptotic}
 h(m)=\lg\lg m+O(\sqrt{\log\log m})\qquad(m\to\infty),
\end{equation}
and in particular $h(m)/\lg\lg m\to1$.
Here additive bounded errors caused by rounding or changing $m$ to $m+1$
are absorbed in the stated error. The lower bound is
$h(m)\ge1+\lg\lg m$ from \eqref{eq:max}.

This conclusion concerns depth, not expression size. Reusing powers of
two in a DAG makes a convenient certificate, but reuse is not needed for
the mathematical result: unfold the DAG into a tree, preserving height.
Every leaf of the unfolded expression is still $1$. No large constant is
introduced for free.

\subsection{Executable realization}
The class \texttt{ExpressionBuilder} implements the proof directly.
It constructs the nonzero block coefficients recursively, builds powers
of two by balanced products of copies of $1+1$, multiplies the terms, and
adds them in a balanced tree. A node table records only three operations:
\texttt{one}, \texttt{add}, and \texttt{mul}.

The supplied independent checker evaluates the node table in topological
order and recomputes every node height. The included 256-bit witness has
165 serialized nodes and height 14; its parameter is $t=8$, for which the
proved bound is $8+4+2=14$. Tests also exercise arbitrary-looking,
all-ones, and single-bit inputs up to 16,384 bits. These are implementation
checks of a construction already proved for all inputs, not empirical
substitutes for Theorem~\ref{thm:radix}.

\section{Proof of the explicit counting bound}\label{sec:sharp-growth}
We now prove the first part of Theorem~\ref{thm:main}.
Fix $n\ge16$ and abbreviate
\[
 t=\ceil{\lg(2n)},\qquad r=r(t),\qquad k=t+r+2,
 \qquad R=2^{n-k}.
\]
The inequality $r\le t$ gives
\[
 k\le2t+2\le2\lg n+6\le n\qquad(n\ge16),
\]
so $R$ is a positive integer power of two.

Take $X=2^{2n}$. As in Section~\ref{sec:first-growth}, every prime at most
$X$ is strictly below $2^{2^t}$. Theorem~\ref{thm:radix} puts all these
primes in $S_k$, and Lemma~\ref{lem:chebyshev} supplies
\[
 P\ge\frac{2^{2n}}{4n\log2}.
\]
Since $2^k\ge2^{t+2}\ge8n$, we have
\[
 \frac{P}{R}\ge\frac{2^{n+k}}{4n\log2}\ge2^n.
\]
Prime-multiset amplification therefore yields
\begin{equation}\label{eq:amplified-final}
 \lg a_n\ge R\lg(P/R)\ge nR=n2^{n-k}.
\end{equation}
Using $2^t<4n$ in the final expression gives
\[
 n2^{n-k}
 =\frac{n2^n}{2^t2^{r+2}}
 >2^{n-r-4}.
\]
Weakening the strict inequality proves the lower half of
\eqref{eq:explicit-main}. The upper half is Proposition~\ref{prop:maximum}.

Finally, \eqref{eq:triangular} gives
\[
 r_n=\sqrt{2\lg n}+O(1).
\]
Taking $\lg$ twice in \eqref{eq:explicit-main} proves
\eqref{eq:doublelog}. Since $\sqrt{\log n}=o(\log n)$, taking the natural
logarithm once gives \eqref{eq:subpoly}.
This establishes the asserted global growth without any sum--product
conjecture or prime-distribution hypothesis.

\begin{center}
\begin{tabular}{r r r r r r}
\toprule
$n$ & $t_n$ & $r_n$ & $k$ & lower bound for $\lg\lg a_n$
 & upper bound\\
\midrule
16 & 5 & 3 & 10 & 9 & 15\\
32 & 6 & 3 & 11 & 25 & 31\\
64 & 7 & 4 & 13 & 56 & 63\\
128 & 8 & 4 & 14 & 120 & 127\\
256 & 9 & 4 & 15 & 248 & 255\\
1024 & 11 & 5 & 18 & 1015 & 1023\\
10000 & 15 & 5 & 22 & 9991 & 9999\\
1000000 & 21 & 6 & 29 & 999990 & 999999\\
\bottomrule
\end{tabular}
\end{center}
These rows evaluate the theorem, not the enormous sequence terms.
For example, the row $n=16$ asserts $a_{16}\ge2^{512}$, not that $a_{16}$
has been enumerated.

\section{One-step growth through a divisor bound}\label{sec:local}
The global estimate alone does not force
$\log a_{n+1}/\log a_n\to2$: broad upper and lower bounds do not rule out
irregular local behavior. We control that behavior separately by bounding
multiplicative collisions.

Let $\tau(z)$ denote the number of positive divisors of $z$.
\begin{lemma}[Uniform divisor upper bound]\label{lem:divisor}
As $N\to\infty$,
\begin{equation}\label{eq:divisor}
 \log\max_{1\le z\le N}\tau(z)
 \le(\log2+o(1))\frac{\log N}{\log\log N}.
\end{equation}
\end{lemma}
\begin{proof}
Put $L=\log N$ and $y=L/(\log L)^3$; take $N$ sufficiently large that
$y>1$. Factor $z=\prod p^{e_p}$. For primes $p\le y$, there are at most
$y$ possible primes, and each exponent is at most $L/\log2$. Their total
contribution to $\log\tau(z)=\sum_p\log(e_p+1)$ is at most
\[
 y\log(1+L/\log2)
 =O\!\left(\frac{L}{(\log L)^2}\right).
\]
For primes $p>y$, the inequality $e+1\le2^e$ gives
\[
 \sum_{p>y}\log(e_p+1)
 \le(\log2)\sum_{p>y}e_p
 \le\frac{(\log2)L}{\log y}.
\]
The last step uses $\sum e_p\log p=\log z\le L$.
Since $\log y=\log L-3\log\log L$, this last expression is
$(\log2+o(1))L/\log L$. Both bounds are uniform in $z\le N$,
proving \eqref{eq:divisor}.
\end{proof}

\begin{proposition}\label{prop:local}
Writing $b_n=\log a_n$, one has
\begin{equation}\label{eq:local-b}
 2b_n-(\log2+o(1))\frac{2^n}{n}
 \le b_{n+1}\le2b_n.
\end{equation}
\end{proposition}
\begin{proof}
Consider all $a_n^2$ \emph{ordered} pairs in $S_n\times S_n$ and map
$(x,y)$ to $xy$. For a fixed product $z$, at most $\tau(z)$ ordered pairs
of positive integers can yield $z$, since the first operand is a divisor
and determines the second. All these products are at most
\[
 M_n^2=M_{n+1}=2^{2^n}.
\]
Hence, with $D_n=\max_{z\le2^{2^n}}\tau(z)$,
\[
 a_{n+1}\ge\card{S_nS_n}\ge\frac{a_n^2}{D_n}.
\]
Apply Lemma~\ref{lem:divisor} at $N=2^{2^n}$. Since
$\log N=(\log2)2^n$ and $\log\log N=n\log2+O(1)$,
\[
 \log D_n\le(\log2+o(1))\frac{2^n}{n}.
\]
Taking logarithms proves the lower bound. The upper bound is
\eqref{eq:square}.
\end{proof}

To finish Theorem~\ref{thm:main}, define
$\varepsilon_n=2-b_{n+1}/b_n$ for $n\ge1$.
Proposition~\ref{prop:local} and \eqref{eq:explicit-main} give
\begin{align*}
 0\le\varepsilon_n
 &\le\frac{(\log2+o(1))2^n}{n\,b_n}\\
 &\le\frac{(16+o(1))2^{r_n}}{n}
 =n^{-1+o(1)}.
\end{align*}
Here $2^{r_n}=\exp(O(\sqrt{\log n}))=n^{o(1)}$.
This proves \eqref{eq:local-main}, with all collision losses accounted for.

\section{The precise computational cost}\label{sec:cost}
\subsection{Candidate counts and unit-cost arithmetic}
If $m=a_{j-1}$, one step has $m(m+1)/2$ unordered pairs and two mathematical
candidate values per pair. Thus the exact candidate count in a cold run
through stage $n$ is
\begin{equation}\label{eq:candidates}
 Q_n=\sum_{j=1}^{n}a_{j-1}(a_{j-1}+1).
\end{equation}
This counts evaluated sum/product candidates, not distinct outputs.
The original program runs the pair iterator twice; the one-pass variant
runs it once and computes both outputs inside the loop. The internal
identity operations used by the built-in \texttt{sum} and \texttt{prod}
only affect constant factors in this analysis.

By Proposition~\ref{prop:cardbounds}, successive cardinalities at least
double. For $m=a_{n-1}$,
\[
 m(m+1)\le Q_n\le\frac43m^2+2m,
\]
so $Q_n=\Theta(a_{n-1}^2)$. In a unit-cost arithmetic model with bounded
amortized set-operation cost, this is also the running-time order.
The assertion about $Q_n$ itself is deterministic and does not need a
hash-table assumption.

For $n=7$, the input size $a_6=67144$ gives
\[
 \binom{67145}{2}=2\,254\,191\,940
\]
unordered pairs, or $4\,508\,383\,880$ candidate values. The output is
reported to contain $357\,306\,081$ distinct integers. This jump explains
why verifying through $n=6$ is modest while casually rerunning the next
stage is not. We did not run that next stage.

\subsection{Bit operations and memory}
The largest integer in $S_n$ has
\[
 B_n=2^{n-1}+1\qquad(n\ge1)
\]
binary digits. Treating its arithmetic as permanently constant-time would
therefore be inappropriate for an asymptotic bit-complexity claim.
Let $\mathsf M(B)$ be an upper bound for multiplying two $B$-bit integers.
Under the usual bounded-average-probe hash-table model, an upper bound for
a complete run is
\begin{equation}\label{eq:bitcost}
 O\!\left(a_{n-1}^2\bigl(\mathsf M(B_{n-1})+B_n\bigr)\right).
\end{equation}
The $B_n$ term covers additions, hashing, and integer comparisons per
bounded number of probes. This is a stated computational model, not a
worst-case guarantee about every hash pattern in a particular interpreter.
Hash collisions affect time but not the mathematical equality semantics
of the set.

Even with every stage cached,
\[
 \sum_{j=0}^{n}a_j<2a_n.
\]
The old stage, the new stage, and the original program's temporary sum,
product, and union sets therefore require $\Theta(a_n)$ stored entries
up to fixed multiplicative factors, not $\Theta(na_n)$.
A bit-storage upper bound is $O(a_n2^n)$. An explicit, uncompressed
collection of $a_n$ distinct positive integers requires
$\Omega(a_n\log a_n)$ bits: at least half of them are at least $a_n/2$.
The ratio between these two bounds is only $n^{o(1)}$ by
\eqref{eq:subpoly}. Object headers and table slack affect practical memory
constants, which are not estimated here.

\subsection{An output-sensitive consequence}
The local growth theorem implies
\[
 \frac{2\log a_{n-1}}{\log a_n}\longrightarrow1.
\]
Therefore the candidate count satisfies
\begin{equation}\label{eq:outputcost}
 Q_n=a_n^{1+o(1)}.
\end{equation}
This is logarithmic-scale near-linearity in the number of output integers,
not near-linearity in $n$. The computation remains doubly exponential in
$n$.

Using even schoolbook multiplication, the arithmetic factors in
\eqref{eq:bitcost} are $2^{O(n)}$; their logarithms are negligible compared
with $\log a_n$. Thus the same $a_n^{1+o(1)}$ conclusion holds for bit time
under the stated hash-table model. A deterministic alternative is to
sort and deduplicate the integer candidates. Its sorting factor and
integer-comparison cost are also $a_n^{o(1)}$, so it attains the same
logarithmic output exponent without a hashing assumption, though its
memory layout and constants differ. No implementation of that large-scale
sorting variant is claimed in the accompanying package.

\section{Recomputed data and verification}\label{sec:verification}
The executable tests compare the one-pass and cached enumerators, a
separate two-comprehension implementation, ordered Cartesian-pair
enumeration through $n=5$, and the value-indexed recurrence
\eqref{eq:h-dp} for every $m\le2000$. These checks exercise distinct ways
of organizing the computation.

Let $L_n$ be the largest integer with $[1,L_n]\subseteq S_n$, and let $P_n$
be the number of primes in $S_n$. The following data were computed in the
supplied run. A dash means that a quantity was not recomputed.
\begin{center}
\begin{tabular}{r r r r r}
\toprule
$n$ & $a_n$ & $L_n$ & $P_n$ & $\beta_n=2^{-n}\lg a_n$\\
\midrule
0 & 1 & 1 & 0 & 0.0000000000\\
1 & 2 & 2 & 1 & 0.5000000000\\
2 & 4 & 4 & 2 & 0.5000000000\\
3 & 11 & 9 & 4 & 0.4324289523\\
4 & 52 & 25 & 9 & 0.3562774824\\
5 & 678 & 177 & 56 & 0.2939106707\\
6 & 67144 & 7390 & 2033 & 0.2505464198\\
7 & $357306081^{\ast}$ & --- & --- & $0.2219733221^{\ast}$\\
\bottomrule
\end{tabular}
\end{center}
\noindent $\ast$ Reported OEIS term; not independently enumerated here.
The small initial gaps are already visible at $n=3$:
\[
 S_3=\{1,2,3,4,5,6,7,8,9,12,16\}.
\]
For $S_6$, the sumset has $21676$ elements, the product set has $55883$,
and their intersection has $10415$, giving
\[
 21676+55883-10415=67144.
\]
These values refer to the sumset and product set formed \emph{from $S_5$}.

The prime counts do not require sieving all the way to $M_n$.
A prime $p\in S_n$ has a minimum-height expression whose root cannot be a
proper product. Its root is a sum of two elements of $S_{n-1}$, so
$p\le2M_{n-1}$. For $n=6$ this is $131072$, despite
$M_6=4294967296$. The data generator sieves only this smaller interval.

All 16 tests in the final test suite pass. They include input validation
before deduplication, immutable-cache behavior, exact maxima, nesting,
cardinality inequalities, power-of-two depths, the triangular recurrence,
and independently evaluated expression certificates. The supplied run used
Python 3.13.5; the implementation targets Python 3.9 or later and has no
third-party dependencies. The runtime and platform are recorded in
\texttt{verification/metadata.json}. The tests are finite validation;
the all-$n$ conclusions come from the proofs above.

\section{What is established, and what remains to refine}\label{sec:scope}
The central established result is not merely that $a_n$ is ``very large.''
Its double logarithm has slope exactly one in base two, with an explicit
one-sided error of order $\sqrt{\log n}$; successive ordinary logarithms
have ratio tending to two. The proof is constructive and every
number-theoretic estimate used in it has been proved in this article.

There is also an exact way to express the remaining normalized constant.
Define
\[
 d_n=\lg\left(\frac{a_n^2}{a_{n+1}}\right)\ge0.
\]
Then
\[
 \beta_n-\beta_{n+1}=\frac{d_n}{2^{n+1}},\qquad
 \theta=\frac12-\sum_{n=1}^{\infty}\frac{d_n}{2^{n+1}}.
\]
The series converges because its partial sums telescope and are bounded
above by $1/2$. Its terms measure logarithmic collision loss relative to
perfect squaring.

The present estimates do not decide whether that series exhausts $1/2$.
If $\theta>0$, then $\log a_n\sim(\theta\log2)2^n$.
If $\theta=0$, the proved lower bound still forces an unusually slow
vanishing: for every $\delta>0$, eventually
$2^{-n}\log a_n\ge n^{-\delta}$. Both possibilities are consistent with
the results proved here. The first eight terms alone cannot choose
between them, and their normalized logarithms are upper bounds for
$\theta$, not estimates with a certified error bar.

A route to a finer result would require tighter information about the
multiplicative collision fibers of the actual sets $S_n$, or a more
efficient source of multiplicatively independent constructible integers.
The uniform divisor bound controls all integers up to $M_{n+1}$, including
exceptional integers that need not be representative of these sets.
Improving that specific overestimate is a concrete mathematical issue;
it is not addressed by additional memoization or by counting expression
syntax.

\appendix
\section{Contents and reproduction instructions}\label{app:reproduce}
The archive includes the following components.
\begin{description}[style=nextline]
\item[\texttt{article.tex}, \texttt{article.pdf}]
The complete mathematical article and its typeset version.
\item[\texttt{code/a352969.py}]
Validated cached and iterative enumerators, a proper-decomposition depth
oracle, exact bound-parameter routines, the constructive expression
builder, and the independent certificate checker.
\item[\texttt{code/test\_a352969.py}, \texttt{code/reproduce.py}]
The regression tests and data-generation script.
\item[\texttt{data/}]
Statistics in CSV, all sets $S_0,\ldots,S_6$ as compressed JSON,
explicit-bound examples, and the expression witness.
\item[\texttt{verification/}]
Test output, data-generation output, and execution metadata.
\end{description}

From the archive root, reproduce the data with
\begin{lstlisting}[language={}]
python code/reproduce.py
python -m unittest discover -s code -p "test_*.py" -v
\end{lstlisting}
To compile the article, run \texttt{python build.py}, or use
\begin{lstlisting}[language={}]
latexmk -pdf -interaction=nonstopmode -halt-on-error article.tex
\end{lstlisting}
\begin{minipage}{\linewidth}
The default reproduction stops at $n=6$. The general enumerator accepts a
\texttt{pair\_budget} keyword so a caller can reject an unexpectedly large
stage before allocating its output. For example:
\begin{lstlisting}
from a352969 import reachable, counting_bound_parameters

assert len(reachable(6, pair_budget=1_000_000)) == 67144
print(counting_bound_parameters(1_000_000))
\end{lstlisting}
\end{minipage}
The parameter routine uses exact integer arithmetic and does not try to
allocate $2^{2^n}$. Its returned exponents specify a theorem-backed bound,
not a numerical enumeration.

\section{Source provenance}
The OEIS page was retrieved on September 20, 2026. Its text-format record
identifies revision 30, dated April 26, 2022. The mathematical definition,
the listed terms, the program attribution, and the wrapper-name typo are
source facts. The correctness proof, constructions, asymptotic estimates,
complexity analysis, and executable package were developed for this note.
No claim of priority over all prior mathematical literature is made.
The original OEIS program is described rather than redistributed verbatim;
the accompanying implementation is independently written.

\begin{thebibliography}{9}
\bibitem{oeis}
V. Reshetnikov, \emph{A352969}, The On-Line Encyclopedia of Integer
Sequences. Python program credited to C. W. Wu, April 15, 2022.
\url{https://oeis.org/A352969}.
Text-format record: \url{https://oeis.org/search?q=id:A352969&fmt=text}.
Accessed September 20, 2026.

\bibitem{itertools}
Python Software Foundation, \emph{itertools---Functions creating iterators
for efficient looping}, documentation of
\texttt{combinations\_with\_replacement}.
\url{https://docs.python.org/3/library/itertools.html#itertools.combinations_with_replacement}.
Accessed September 20, 2026.

\bibitem{functools}
Python Software Foundation, \emph{functools---Higher-order functions and
operations on callable objects}, documentation of \texttt{lru\_cache}.
\url{https://docs.python.org/3/library/functools.html#functools.lru_cache}.
Accessed September 20, 2026.

\bibitem{types}
Python Software Foundation, \emph{Built-in Types}, integer arithmetic and
set/frozenset semantics.
\url{https://docs.python.org/3/library/stdtypes.html}.
Accessed September 20, 2026.

\bibitem{math}
Python Software Foundation, \emph{math---Mathematical functions},
documentation of \texttt{prod}.
\url{https://docs.python.org/3/library/math.html#math.prod}.
Accessed September 20, 2026.

\bibitem{complexity}
The OEIS Foundation, \emph{A005245}, integer complexity measured by the
minimum number of ones.
\url{https://oeis.org/A005245}.
Accessed September 20, 2026.

\bibitem{depth}
The OEIS Foundation, \emph{A210659}, depth among representations having
minimum operand complexity.
\url{https://oeis.org/A210659}.
Accessed September 20, 2026.
\end{thebibliography}
\end{document}
